% SPDX-License-Identifier: AGPL-3.0-or-later
% Commercial license available
% © Concepts 1996–2026 Miroslav Šotek. All rights reserved.
% © Code 2020–2026 Miroslav Šotek. All rights reserved.
% ORCID: 0009-0009-3560-0851
% Contact: www.anulum.li | protoscience@anulum.li
% SCPN-FUSION-CORE — same-case free-boundary comparison with FreeGSNKE.
\documentclass[11pt]{article}

\usepackage[utf8]{inputenc}
\usepackage[T1]{fontenc}
\usepackage{amsmath}
\usepackage{booktabs}
\usepackage{microtype}
\usepackage{natbib}
\usepackage[colorlinks=true,linkcolor=blue,citecolor=blue,urlcolor=blue]{hyperref}
\usepackage[margin=1in]{geometry}
\bibliographystyle{unsrtnat}

\title{Same-Case Comparison of Two Free-Boundary Grad--Shafranov Codes on
FreeGSNKE's Public Machine Descriptions:\\
Forward Parity and Inverse Coil-Current Fits}
\author{Miroslav \v{S}otek\\
\small ANULUM CH\&LI, Marbach SG, Switzerland, \url{https://www.anulum.li}\\
\small ORCID: \href{https://orcid.org/0009-0009-3560-0851}{0009-0009-3560-0851}\\
\small \href{mailto:protoscience@anulum.li}{protoscience@anulum.li}}
\date{Preprint --- 16 November 2026}

\begin{document}
\maketitle

\begin{abstract}
We compare the predictive free-boundary Grad--Shafranov solver of
SCPN-Fusion-Core with FreeGSNKE 3.0.1 on the three static inverse problems that
FreeGSNKE ships, for synthetic MAST-U-like, SPARC-like and ITER-like machine
descriptions.  A bridge maps FreeGSNKE's coils, profiles and flux convention
onto SCPN's: the vacuum fluxes agree to rounding, and the sampled profile
reproduces FreeGSNKE's current density to a relative $L_2$ error below
$10^{-7}$.  In the forward direction, started from FreeGSNKE's flux, SCPN reaches
its own fixed point within a root-mean-square normalised-flux difference of
$3.5 \times 10^{-3}$, $4.6 \times 10^{-4}$ and $6.4 \times 10^{-4}$ inside the
limiter; imposing FreeGSNKE's plasma support lowers these by factors of 3 to 9.
From the vacuum field SCPN reaches the same solution on the SPARC- and ITER-like
cases but finds no magnetic axis on the MAST-U-like case.

In the inverse direction SCPN fits its coil currents to FreeGSNKE's isoflux and
null-point targets, and each code's currents are scored by both codes.  On
FreeGSNKE's flux the two evaluators agree to 0.3\%.  Started from FreeGSNKE's
currents, SCPN meets the MAST-U- and SPARC-like targets to below $5 \times 10^{-8}$ with current changes of 1.4\% and
1.0\%.  FreeGSNKE scores those currents above its own solution on the
MAST-U-like case and 4.5 times below it on the SPARC-like case.  From a
perturbed start SCPN stops early on the MAST-U-like case, and FreeGSNKE does
not converge at its currents, so that comparison is incomplete.  Neither code
meets the overdetermined ITER-like targets.  The SPARC-like Jacobian is
rank-deficient, and two current sets 12.7\% apart both meet the targets.  Codes
should therefore be compared by their target residuals under both evaluators,
not by their coil currents.  All cases are synthetic: no measured discharge
enters, and no speed, control or facility claim is made.
\end{abstract}

\noindent\textbf{Keywords:} free-boundary equilibrium; Grad--Shafranov;
code comparison; inverse equilibrium; coil-current fit; FreeGSNKE

\section{Introduction}

Free-boundary equilibrium codes solve the Grad--Shafranov equation
\citep{GradRubin1958,Shafranov1958} for the poloidal flux of a plasma and its
external coils.  When two such codes are compared on configurations that each
code set up for itself, the differences mix the numerical methods with the
inputs: the machine description, the profile parametrisation, the grid and the
conventions for flux and current.  A same-case comparison fixes every input and
leaves only the methods.

FreeGSNKE \citep{Amorisco2024} is an open-source free-boundary solver.  It
ships machine descriptions and static inverse problems for three
configurations, and its static forward solver has been validated against
EFIT++ reconstructions of MAST-U discharges together with the Fiesta code
\citep{Pentland2025}.  That study treats the static forward solvers only; it
omits inverse and reconstruction solvers and describes a review of them as
beyond its scope \citep{Pentland2025}.

This paper compares the predictive free-boundary solver of SCPN-Fusion-Core
with FreeGSNKE 3.0.1 on those three shipped problems, in two directions.  In
the forward direction, SCPN solves the equilibrium from FreeGSNKE's coil
currents and profile; the comparison is the normalised flux inside the limiter.
In the inverse direction, SCPN fits its own coil currents to FreeGSNKE's
isoflux and null-point targets, and each code's currents are scored by both
codes.  The contribution is the bridge that makes the two codes solve the same
problem, the measured agreement and disagreement in both directions, and the
failures that the comparison exposed.

\section{Codes and cases}

\subsection{FreeGSNKE and its public problems}

FreeGSNKE inherits the Picard-based static inverse solver of FreeGS
\citep{Dudson2024} and adds a Newton--Krylov static forward solver
\citep{Amorisco2024,Pentland2025}.  We use version 3.0.1 at source commit
\texttt{00b79a08fbb6}, with FreeGS4E 0.13.1.  Its repository ships three
synthetic machine descriptions; according to its README, the MAST-U-like
description comes from UKAEA, the SPARC-like one from the public SPARC
repository, and the ITER-like one from FUSE.jl.  Each case is the static inverse
problem of an upstream test or example (Table~\ref{tab:cases}).  FreeGSNKE
solves it, and a forward refinement at the resulting currents converges to a
relative flux change below $10^{-9}$.  That refined equilibrium is the reference
for both directions.

\begin{table}[ht]
\centering
\caption{The three public cases.  Grid sizes are $n_R \times n_Z$.  Active
filaments inside the computational rectangle require the coil field to be kept
separate from the plasma field (Section~\ref{sec:bridge}).}
\label{tab:cases}
\begin{tabular}{lrrrrr}
\toprule
Case & $I_p$ (MA) & Grid & Circuits & Filaments (inside) & Profile class \\
\midrule
MAST-U-like & 0.6 & $65 \times 129$ & 12 & 876 (876) & \texttt{ConstrainPaxisIp} \\
SPARC-like & 8.7 & $129 \times 129$ & 11 & 1380 (150) & \texttt{ConstrainPaxisIp} \\
ITER-like & 15.0 & $129 \times 129$ & 10 & 20 (8) & \texttt{Fiesta\_Topeol} \\
\bottomrule
\end{tabular}
\end{table}

\subsection{The SCPN predictive solver}

SCPN-Fusion-Core's predictive solver finds the equilibrium from the coil
currents, the profile functions $p'(\psi)$ and $FF'(\psi)$, and the plasma
current.  The flux on the boundary of the computational rectangle is the coil
flux plus the plasma's own flux there, written as a precomputed Green's-function
response to the toroidal current density of every interior cell.  The current
density is rescaled each iteration to the target plasma current.  The magnetic
axis and the separatrix flux follow from smooth O- and X-point estimates of the
current iterate, so the last closed flux surface is an output.  Plain Picard
iteration on this coupled map is unstable, so the fixed point is found by
Anderson mixing \citep{Anderson1965,WalkerNi2011}.  A cold start ramps the
plasma current from the vacuum field; a warm start begins from a given flux.

\subsection{The bridge between the codes}
\label{sec:bridge}

Four audits make the two codes solve the same problem.  First, FreeGSNKE's
resolved circuit currents are mapped onto SCPN's filament Green's functions,
and the two vacuum fluxes are compared inside the limiter.  Second, the
converged FreeGSNKE profile is sampled on 8193 monotonic $\psi_N$ knots onto
SCPN's $p'$ and $FF'$ source surface, with $\psi_N = (\psi -
\psi_\mathrm{axis})/(\psi_\mathrm{boundary} - \psi_\mathrm{axis})$; the toroidal
current density built by SCPN from those samples is compared with FreeGSNKE's
on FreeGSNKE's plasma support.  Of five candidate flux conventions (identity,
scaled by $\pm 2\pi$, by $1/2\pi$ and by $-1$), only the identity reproduces the
source, at a relative $L_2$ error below $10^{-7}$, while the others miss by
0.84 or more; both codes therefore use flux per radian with the same sign.
Third, a constant flux offset leaves the source unchanged, as it must.  Fourth,
active filaments lie inside the computational rectangle in all three cases, so
SCPN forms the total flux as the sum of the coil flux and the plasma flux and
keeps the two fields separate.

\section{Forward parity}

Table~\ref{tab:forward} lists the forward comparison.  The vacuum fluxes agree
to rounding.  The sampled source reproduces FreeGSNKE's current density to a
relative $L_2$ error of at most $9.2 \times 10^{-8}$.  Started from FreeGSNKE's
flux, SCPN converges to its own fixed point, and the normalised flux of the two
equilibria differs inside the limiter by a root-mean-square $\psi_N$ of
$3.5 \times 10^{-3}$ on MAST-U-like and $4.6$ and $6.4 \times 10^{-4}$ on the
SPARC- and ITER-like cases.  The magnetic axes agree to within 1.4~mm and the
X-points to within 2.1~mm.

\begin{table}[ht]
\centering
\caption{Forward parity.  Vacuum: largest absolute vacuum-flux difference
inside the limiter.  Source: relative $L_2$ difference of the toroidal current
density.  Warm: SCPN's own smooth topology.  Frozen: FreeGSNKE's plasma
support imposed.  Axis and X-point: position differences of SCPN's warm
equilibrium.}
\label{tab:forward}
\begin{tabular}{lrrrrrr}
\toprule
 & Vacuum & Source & \multicolumn{2}{c}{$\psi_N$ RMSE} & Axis & X-point \\
Case & (Wb\,rad$^{-1}$) & rel.\ $L_2$ & warm & frozen & (mm) & (mm) \\
\midrule
MAST-U-like & $1.1 \times 10^{-16}$ & $3.5 \times 10^{-8}$ & $3.54 \times 10^{-3}$ & $3.78 \times 10^{-4}$ & 1.35 & 2.11 \\
SPARC-like & $1.8 \times 10^{-14}$ & $9.2 \times 10^{-8}$ & $4.59 \times 10^{-4}$ & $1.52 \times 10^{-4}$ & 0.12 & 1.19 \\
ITER-like & $7.7 \times 10^{-13}$ & $3.0 \times 10^{-8}$ & $6.42 \times 10^{-4}$ & $1.22 \times 10^{-4}$ & 1.26 & 1.50 \\
\bottomrule
\end{tabular}
\end{table}

The frozen-support control separates the two sources of that difference.  When
SCPN is given FreeGSNKE's plasma support instead of finding its own, the
$\psi_N$ difference falls by a factor of 9.4 on MAST-U-like, 3.0 on SPARC-like
and 5.3 on ITER-like.  Most of the warm difference therefore comes from where
each code places the plasma boundary, not from the source or the field solve.
The thresholds of the comparison were fixed before the runs: a $\psi_N$ RMSE of
at most 0.05; axis and X-point differences of at most 35~mm; a relative
support-current difference of at most 0.1; a relative nonlinear residual of at
most 0.05; a relative plasma-current error of at most $10^{-6}$; and relative
errors of the implicit gradients, checked by finite differences, of at most
0.05 for the coil currents and 0.01 for the profile.  All three cases pass every
threshold.

The cold start is a robustness test of SCPN alone.  From the vacuum field, the
SPARC-like and ITER-like cases reach the warm-start fixed point, after 153 and
270 iterations, with the same $\psi_N$ difference to three significant digits.
The MAST-U-like case does not: after 600 iterations it has no magnetic axis
inside the limiter.  The Grad--Shafranov problem can have several solutions
\citep{HamFarrell2024}, and from the vacuum field SCPN does not reach the
diverted solution that the warm start finds.
The forward agreement on MAST-U-like is therefore a property of the warm start.

\section{Inverse coil-current fits}

\subsection{Targets and method}

Each upstream problem names isoflux points, which must lie on one flux surface,
and two null points, where the poloidal field must vanish
(Table~\ref{tab:targets}).  The residual has three blocks: the flux
differences of every pair of isoflux points, and $B_R$ and $B_Z$ at the null
points.  We report the $L_2$ norm of each block and their root sum of squares
as the total, in the mixed units FreeGSNKE uses (Wb\,rad$^{-1}$ and T).  With
$n$ isoflux points the pairs fix only $n-1$ independent flux differences, so
the three problems pose 11, 9 and 27 independent conditions on 11, 11 and 10
controlled circuits.

\begin{table}[ht]
\centering
\caption{Inverse targets.  Independent conditions: $n-1$ flux differences plus
two field components at each null point.}
\label{tab:targets}
\begin{tabular}{lrrrr}
\toprule
Case & Isoflux points & Null points & Independent conditions & Controlled circuits \\
\midrule
MAST-U-like & 8 & 2 & 11 & 11 \\
SPARC-like & 6 & 2 & 9 & 11 \\
ITER-like & 24 & 2 & 27 & 10 \\
\bottomrule
\end{tabular}
\end{table}

SCPN fits the same controlled circuits by a Levenberg--Marquardt iteration
\citep{Levenberg1944,Marquardt1963} with Marquardt scaling, adaptive damping, the
upstream Tikhonov weights on the step and a projection onto the upstream current
limits.  The flux at a target point is the exact filament coil flux plus a Keys
cubic-convolution interpolation \citep{Keys1981} of the plasma flux.  The
Jacobian is taken by central differences through the full free-boundary solve,
and every forward solve must converge to $10^{-9}$ or the step is rejected.  The
methods differ from FreeGSNKE's in three ways: FreeGSNKE linearises the targets
with the plasma frozen, penalises limit violations through slack variables
instead of projecting, and interpolates the plasma flux with a bicubic spline.
Both apply the Tikhonov weight to the current step.  Each fit starts either
from FreeGSNKE's currents or from a seeded $\pm 20\%$ perturbation of them,
reached by continuation from FreeGSNKE's currents when the perturbed currents do
not converge directly.  A fit stops when no step lowers the loss or after 30
outer iterations.  FreeGSNKE is then refined at SCPN's fitted currents, so each
code's currents are scored by both codes.

\subsection{Results}

Table~\ref{tab:inverse} lists the losses.  The two evaluators agree: applied to
FreeGSNKE's own flux map, SCPN's evaluator reproduces FreeGSNKE's total loss to
0.3\% in all three cases.  The $B_R$ block, a derivative of the flux, differs more: by up
to 26\% on the SPARC-like case.  The two codes interpolate the plasma flux
differently.  The totals that follow therefore differ mainly through the two
equilibria, not through the evaluators.

\begin{table}[ht]
\centering
\caption{Inverse losses (total $L_2$).  At FreeGSNKE's currents: FreeGSNKE's own
loss, SCPN's evaluator on FreeGSNKE's flux map, and SCPN's own equilibrium.  At
SCPN's fitted currents: SCPN's loss and FreeGSNKE's.  Change: relative $L_2$
change of the controlled currents from FreeGSNKE's.  Stop: no descent step (n),
the 30-iteration limit (m), or a forward solve in the Jacobian that did not
converge (j).  Withheld: FreeGSNKE's refinement did not converge.}
\label{tab:inverse}
\footnotesize
\setlength{\tabcolsep}{3pt}
\begin{tabular}{llrrrrrrl}
\toprule
 & & \multicolumn{3}{c}{At FreeGSNKE's currents} & \multicolumn{2}{c}{At SCPN's currents} & & \\
\cmidrule(lr){3-5}\cmidrule(lr){6-7}
Case & Start & FreeGSNKE & SCPN eval. & SCPN & SCPN & FreeGSNKE & Change & Stop \\
\midrule
MAST-U-like & FreeGSNKE & $1.54 \times 10^{-3}$ & $1.55 \times 10^{-3}$ & $2.24 \times 10^{-3}$ & $1.92 \times 10^{-8}$ & $2.11 \times 10^{-3}$ & 1.40\% & n \\
 & perturbed & & & & $1.96 \times 10^{-2}$ & withheld & 64.2\% & j \\
SPARC-like & FreeGSNKE & $3.49 \times 10^{-2}$ & $3.48 \times 10^{-2}$ & $3.11 \times 10^{-2}$ & $4.92 \times 10^{-8}$ & $7.78 \times 10^{-3}$ & 0.97\% & n \\
 & perturbed & & & & $1.21 \times 10^{-6}$ & $7.75 \times 10^{-3}$ & 12.7\% & n \\
ITER-like & FreeGSNKE & 1.71 & 1.71 & 1.81 & 1.39 & 1.39 & 1.26\% & m \\
 & perturbed & & & & 1.41 & 1.41 & 1.21\% & m \\
\bottomrule
\end{tabular}
\end{table}

At FreeGSNKE's currents, SCPN's own equilibrium scores $2.24 \times 10^{-3}$ on
MAST-U-like, 1.45 times FreeGSNKE's loss, and 0.89 and 1.06 times FreeGSNKE's
loss on the SPARC- and ITER-like cases.  This is the model difference between
the codes, measured in target residuals.

Started from FreeGSNKE's currents, SCPN meets the MAST-U-like and SPARC-like
targets to $1.9 \times 10^{-8}$ and $4.9 \times 10^{-8}$ by moving the currents by
1.40\% and 0.97\%.  FreeGSNKE scores those currents at $2.11 \times 10^{-3}$ on
MAST-U-like, above its own $1.54 \times 10^{-3}$, and at $7.78 \times 10^{-3}$ on
SPARC-like, 4.5 times below its own $3.49 \times 10^{-2}$.  SCPN's currents stay inside the current limits, so on the
SPARC-like case the currents that FreeGSNKE's inverse solve returned with the
upstream settings are not a minimiser of its own target loss.  Neither code meets the ITER-like
targets, which overdetermine the ten circuits: SCPN stops at the iteration
limit with 1.39, and FreeGSNKE scores SCPN's currents at the same 1.39, below
the 1.71 of its own solution.

The singular values of SCPN's Jacobian show why currents are a poor basis for
comparison.  On the SPARC-like case six singular values lie between
$2.5 \times 10^{-4}$ and $6.9 \times 10^{-7}$ and the remaining five between
$1.2 \times 10^{-10}$ and $2.2 \times 10^{-21}$ (start from FreeGSNKE's currents);
two of them are zero to rounding, as the nine independent conditions require.  The perturbed start meets the
targets to $1.2 \times 10^{-6}$ with currents 12.7\% away from FreeGSNKE's, and
FreeGSNKE scores both SCPN solutions alike ($7.78$ and
$7.75 \times 10^{-3}$).  On the MAST-U-like case seven singular values lie above
$2 \times 10^{-6}$ and four below $10^{-10}$.  The ITER-like spectrum decays from
$2.2 \times 10^{-4}$ to $1.9 \times 10^{-9}$ without a gap of several orders.
Where the spectrum has such a gap, the targets leave several current
combinations nearly free: on the SPARC-like case two current sets 12.7\% apart
both meet the targets to $1.2 \times 10^{-6}$ or better.  We therefore compare the codes by their target
residuals under both evaluators, not by their currents.

From the perturbed start the MAST-U-like fit fails.  SCPN reaches the
perturbed currents by 387 continuation steps and lowers its loss over seven
iterations to $1.96 \times 10^{-2}$.  It then stops because a forward solve
inside the Jacobian does not converge.  Its currents then differ from
FreeGSNKE's by 64.2\%, and FreeGSNKE's refinement at those currents does not
converge either: its relative change stays at $2.6 \times 10^{-5}$ against the
required $10^{-9}$.  The benchmark withholds that loss.  Its completeness gate,
which requires every FreeGSNKE refinement at SCPN's currents to converge,
fails the case, so the inverse report's status is fail.  The gate was declared
before the runs and is kept.

\section{Limitations}

\begin{itemize}
\item All three machine descriptions are synthetic configurations shipped with
FreeGSNKE.  No measured discharge enters the comparison.
\item Each case uses one grid.  Grid convergence of the differences above is not
shown.
\item The forward agreement is a warm-start result.  SCPN's cold start fails on
the MAST-U-like case.
\item The inverse comparison uses one perturbation seed per case.  On the
ITER-like case neither code meets the targets, so the comparison there is
between residual levels.
\item Run times are recorded only to describe the runs.  The two codes were not
timed on the same problem, tolerance and host, so no speed comparison is made.
\end{itemize}

\section{What this comparison does not show}

It does not validate either code against experiment, and it does not show
fitness for a plasma control system, for facility use or for safety analysis.
FreeGSNKE's inverse solver is the one it inherits from FreeGS; this paper
reports how its solutions score under a second code, not a defect of it.

\section*{Software version and date}

This preprint describes \texttt{scpn-fusion} at commit \texttt{86649a05a15a}
of 16 November 2026, on the main line after release 4.1.0; the comparison code
is not part of that release.  The forward report was produced at commit
\texttt{422e4c0b3d15}.  The six inverse runs ran at commit
\texttt{44852feeedf2}; commit \texttt{86649a05a15a} merged them and tracks the
report.  The solver and benchmark sources named in the
evidence manifest are byte-identical between the inverse run and the commit
above.  Between the two runs no solver source changed; the shared benchmark
helper changed only by naming two constant tables, with the same values.  All runs used one host (Intel Xeon E5-2640, 24 logical CPUs, Python
3.12.3, NumPy 1.26.4, SciPy 1.15.3, JAX 0.7.1).

\section*{Data and code availability}

Code is available at \url{https://github.com/anulum/scpn-fusion-core}.  The
two reports cited here, with their source revisions, Git blob identities and
SHA-256 digests, are in \texttt{evidence/} of this submission directory.  The
generator in that directory checks every quoted value against them.  This is a
venue-neutral preprint and has not been submitted.  No DOI has been reserved
and the manuscript-content licence has not yet been selected by the author.

\section*{Use of AI tools}

Large language model coding agents were used substantively for generating and
debugging the analysis code and for drafting and revising the manuscript text.
The agents were commercially available models, which are not named because no
vendor is involved in or sponsors this work.  They were coordinated through
SYNAPSE CHANNEL (version 0.99.52 at the time of writing;
\url{https://github.com/anulum/synapse-channel}), the open-source
agent-coordination tool of the ANULUM Agentic Coordination, Assurance and
Continuity Systems group.  The author specified the comparison, verified the
code through its tests and every reported number against the committed
reports, and takes full responsibility for the content.

\bibliography{references}
\end{document}
